\section{Commercial instruments and credit eligibility thresholds}
\label{exec13:commercial-iv-rd}

A credit committee changes an eligibility rule at a score threshold. Approval
increases above the threshold, but some eligible borrowers decline the loan,
some ineligible borrowers receive an exception, and a second benefit may
change at the same score. The committee wants the value of extending credit,
not merely the difference in outcomes between score groups. This example
connects instrumental variables (IV) and regression discontinuity (RD), while
showing why selecting a familiar command does not settle identification.
The software must represent the assignment rule, actual action, running
variable, cutoff, observation horizon and uncertainty in the local first
stage. The analyst must defend what the assignment changes besides the action.

Commercial software already provides substantial parts of this workflow.
Stata documents weak-instrument test inversion with disconnected, unbounded
and empty sets. SAS and RATS provide systems and instrument specifications;
EViews provides several distinct instrument diagnostics; SPSS exposes a
two-stage least-squares procedure; and JMP documents equation-specific
model-implied instruments. A proposal that compares only conventional least
squares with its own set-valued interface would miss these alternatives.
The relevant comparison concerns the target, assumptions, numerical boundary,
licensed execution and preservation of the result between applications.

The comparison below uses inspected official technical text for six commercial
families and pinned upstream Stata RD packages. MATLAB, PcGive and LIMDEP
remain candidates with unavailable essential originals in this bounded
review. They are not classified as lacking IV or RD. Documentation identifies
the available operation in its stated edition; it does not establish runtime
performance or a successful integration here. Mutable manuals acquired on
13 October 2026 do not prove their exact bytes existed at the 7 October
literature cutoff. Dated editions and pinned package commits narrow that
uncertainty, and unresolved edition and execution rights remain explicit
\citep{exec13_ivrd_survey}.

\subsection{A threshold jump is a target with assumptions}

Write $S$ for the score, $c$ for the eligibility threshold,
$Z=\mathbf 1\{S\ge c\}$ for eligibility, $A$ for actual credit extension
and $Y$ for subsequent value. Where the relevant limits exist and are
identified by the score support, define
\[
 \Delta_Y=\lim_{s\downarrow c}E[Y\mid S=s]
          -\lim_{s\uparrow c}E[Y\mid S=s],\qquad
 \Delta_A=\lim_{s\downarrow c}E[A\mid S=s]
          -\lim_{s\uparrow c}E[A\mid S=s].
\]
The ratio $\Delta_Y/\Delta_A$ is a local observed contrast when
$\Delta_A\ne0$. Its causal interpretation needs a model. Under continuity
of the relevant potential-outcome and compliance laws, an exclusion
restriction, a nonzero first-stage jump and local monotonicity, a binary
action can give a local complier target. Monotonicity rules out borrowers
whose treatment response moves in the opposite direction to eligibility.
It does not follow from observing a positive aggregate first stage. A
constant structural-effect model instead supplies its own stronger
interpretation. Neither target is automatically the bank-wide effect of
extending credit to every applicant.

The local complier calculation can be seen directly in a stipulated
limiting population whose potential outcomes and treatment responses
have the same distribution on both sides of the cutoff. Under exclusion
and consistency, write $A(1),A(0)$ for eligibility responses and
$Y(1),Y(0)$ for action outcomes. Monotonicity gives
$A(1)-A(0)=\mathbf 1\{C\}$, where $C$ denotes a complier. Consequently
$\Delta_A=P(C)$ and
$\Delta_Y=E[(Y(1)-Y(0))\mathbf 1\{C\}]$; their ratio is
$E[Y(1)-Y(0)\mid C]$ when $P(C)>0$. The common limiting population
is the scientific premise. If eligibility itself changes the outcome
or the two sides approach different potential-outcome laws, this
decomposition does not identify the proposed action effect.

A finite authored construction exposes the exclusion issue without relying
on a fitted estimator. Let action probabilities below and above the threshold
be $.2$ and $.6$, and mean outcomes be USD~90 and USD~110. In one model,
the outcome mean is $80+50E[A\mid Z]$, so the causal coefficient is USD~50
and the reduced-form ratio is also 50. In another, the mean is
$84+30E[A\mid Z]+8Z$. It has exactly the same two action probabilities
and outcome means, but a causal coefficient of USD~30 and a direct
eligibility benefit of USD~8. With an expressly stipulated incremental
cost of USD~40, the implied value of the action is $+10$ or $-10$.
These equal summary moments are not asserted to be equal full distributions.
They demonstrate which additional restriction the ratio uses
\citep{exec13_ivrd_exact}.

More instruments do not automatically validate exclusion. Consider
independent fair signs $Z_1,Z_2,V,U$ and the structural equations
\[
 A=Z_1+2Z_2+V,\qquad
 Y=30A+10(Z_1+2Z_2)+U.
\]
An intervention changing $A$ while holding its causes fixed has coefficient
30. Nevertheless $Y-40A=U-10V$, so both population instrument moments
$E[Z_j(Y-40A)]$ equal zero. Instrument direct effects proportional to
their first-stage contributions escape this population overidentification
comparison. The exact sixteen-state enumeration checks the calculation;
it is not a vendor test run or a claim about a particular finite-sample
rejection probability. First-stage relevance, overidentifying restrictions
and causal instrument validity therefore need separate fields and explanations
\citep{exec13_ivrd_exact}.

\subsection{Stata already provides weak-instrument sets}

The inspected \texttt{ivregress} manual documents two-stage least squares,
limited-information maximum likelihood and generalized method of moments.
The postestimation manual distinguishes endogeneity tests, first-stage
statistics, overidentification tests and \texttt{estat weakrobust}.
For just-identified models the latter defaults to an Anderson--Rubin test;
for overidentified models it defaults to a conditional likelihood-ratio
test, with the documented generalization for a nonhomoskedastic covariance
estimate. The requested null concerns coefficients of the endogenous
regressors. Confidence-set output through \texttt{ci} is restricted to
a single endogenous regressor. A multiple-coefficient hypothesis test
must not be exported as though it were a fully computed joint confidence
region~\citep{exec13_ivrd_stata}.

The manual explicitly describes five confidence-set forms: a finite
interval, a union of finite intervals, a union that can include infinite
intervals, the whole real line and the empty set. Thus preserving
disconnected or unbounded answers is already part of a documented
commercial alternative. The proposed library's general quadratic-set
operation does not by itself establish a statistical advance over these
procedures. Its potential contribution is a common scientific contract
across tools, complete request accounting and a decision layer that
preserves the uncertainty and its construction.

The numerical distinction matters. For some cases Stata computes the
inversion directly; for nonhomoskedastic overidentified CLR inference it
uses a grid. The inspected documentation specifies an initial grid of
500 points, default bounds five conventional standard errors around the
estimated coefficient, adaptive expansion, a default maximum of 20
iterations, and options controlling bounds and adaptation. Where critical
values require simulation, the default is 25,000 draws and a seed can be
supplied. These are documented controls, not measured execution costs.
An adapter should retain the requested and final bounds, grid size,
simulation seed/state, replications, covariance specification and all
interval components~\citep{exec13_ivrd_stata}.

The manual distinguishes a model-rejecting empty Anderson--Rubin set
in an overidentified model from an apparent empty answer produced by
insufficient grid bounds in other cases. This prevents a downstream
agent from reading every empty set as scientific rejection. A simple
authored rectangle makes the danger concrete: $M\in[1,2]$ and
$Q\in[-1/2,1/2]$ imply the ratio relation set
$(-\infty,-2]\cup[2,\infty)$. A grid confined to $[-1,1]$ finds no
accepted points despite a nonempty set. This is an exact geometry example,
not a reproduction of Stata's AR or CLR algorithm. It motivates recording
whether the numerical search actually resolved the declared domain
\citep{exec13_ivrd_exact}.

The official feature pages associate weak-instrument inference with Stata~19
and identify its earlier StataNow availability in the Stata~18 feature
listing. That evidence does not imply that every fixed Stata~18 installation
contains the same command or that the mutable manual is identical across
updates. The deployment manifest must bind the actual installed edition,
update and command help. The inspected licensing page distinguishes an
identified single user from a specified number of concurrent network users.
It does not settle unattended service accounts, embedding, redistribution
or the commercial terms of a proposed MCP service. Those rights and prices
need a relevant license review before a licensed integration trial
\citep{exec13_ivrd_stata,exec13_ivrd_license}.

\subsection{Systems, diagnostics and model-implied instruments are different alternatives}

SAS/ETS~14.1's \texttt{PROC SYSLIN} documentation supplies a useful
systems-oriented comparison. It documents 2SLS, LIML, K-class, MELO,
SUR, 3SLS and full-information alternatives. Its technical discussion
states the instrument correlation and orthogonality requirements. The
\texttt{INSTRUMENTS} statement supplies all listed instruments to each
first-stage regression; an endogenous variable used as a regressor is
declared in \texttt{ENDOGENOUS}. The worked 2SLS example explains the
\texttt{FIRST} option and warns that its ANOVA and $R^2$ need a different
interpretation from ordinary least squares. A consumer therefore needs
the equation system, restrictions, instrument lists and actual estimation
sample, not just a coefficient table. These inspected pages establish
documented operations; they do not settle weak-IV set inversion or RD
support in every SAS edition. The incomplete download of a later manual
is excluded from technical evidence~\citep{exec13_ivrd_sas}.

EViews' instrument-diagnostic documentation distinguishes instrument
summary, subset orthogonality, regressor endogeneity, weak-instrument
diagnostics and a GMM breakpoint test. The summary can expose instruments
dropped for collinearity. Its weak-instrument view reports Cragg--Donald
statistics, available Stock--Yogo critical values and moment-selection
criteria for the stated estimation methods. The manual itself qualifies
the Cragg--Donald statistic as applicable to TSLS and other K-class
estimators while also displaying it for GMM comparison. The proposed
interface should carry that diagnostic qualification. A relevance screen
is neither a confidence set robust to weak identification nor verification
of exclusion. The extracted document's equations contain image elements;
only inspected textual definitions and scope are used here, without
reconstructing unread formulas~\citep{exec13_ivrd_eviews}.

IBM SPSS Statistics~30 documents a 2SLS procedure requiring the Standard
edition or Regression option. It exposes explanatory and instrumental
variable lists, coefficient/covariance output and command syntax. A
regressor omitted from the instrument list is treated as endogenous;
the documentation normally includes the exogenous explanatory variables
in both lists. This is useful for a guided analysis interface, but a
generic menu explanation calling a fitted proxy optimal cannot replace
the identifying assumptions, rank conditions or the appropriate covariance
argument. The inspected procedure page does not establish all-version
weak-IV or RD coverage. Its listed model-based intervals must retain
their own assumptions when exported~\citep{exec13_ivrd_spss}.

RATS~11 provides a programmable alternative. Section~2.5 of its 2025
user's guide specifies the \texttt{INSTRUMENTS} list and \texttt{INST}
option across its regression procedures, with system operations using
the documented GMM route. It writes the orthogonality condition explicitly
and distinguishes the covariance choice through \texttt{ROBUSTERRORS}
and other options. No variable, including the constant, is automatically
included in the instrument list. This small semantic difference can
change a translated model. A generated command file must therefore
declare the constant, included exogenous variables, excluded instruments,
sample window and covariance options explicitly. The source does not
support a claim that an arbitrary RATS program is a valid IV design,
or that every inference extension has been inspected~\citep{exec13_ivrd_rats}.

JMP~19's September 2025 multivariate manual documents model-implied
instrumental-variable 2SLS inside its structural-equation modeling
platform. It estimates equations separately and reports equation-level
Sargan diagnostics and the instrument set for each equation. This is a
valuable alternative when the analyst begins with a supplied structural
model. The instruments are implied by that model; discovering a valid
credit encouragement from an unrestricted business brief is a different
task. The documented local-misspecification advantages are not a theorem
that every structural misspecification or invalid instrument is harmless.
Edition entitlement and an actual scripted export remain to be checked
for the intended installation. A model-implied instrument should be
identified as such in the shared contract~\citep{exec13_ivrd_jmp}.

MATLAB's \texttt{ivx}/\texttt{iv4}, PcGive's instrument-estimation path,
and LIMDEP/NLOGIT remain named acquisition leads. Their essential technical
pages were inaccessible through the bounded official routes used here.
Consequently this comparison makes no technical equivalence, absence,
performance or deployment-rights claim for them. The concrete follow-up
is an edition-matched original manual and license record, followed by
the same instrument, covariance, numerical-set and export comparison.
This omission limits commercial coverage; it does not invalidate the
supported findings for the acquired sources~\citep{exec13_ivrd_survey}.

\subsection{An open RD package on a commercial host has two reuse decisions}

The pinned upstream Stata \texttt{rdrobust} help and source identify
version~11.1.1, dated 1 October 2026. They document local-polynomial
point estimation, robust bias-corrected inference, sharp and fuzzy
designs, distinct estimation and bias bandwidths, covariates and several
variance options. The help distinguishes conventional inference,
bias-corrected estimates with conventional variance, and bias-corrected
estimates with robust variance. A result containing only one coefficient
and standard error loses that distinction. Its MSE and coverage-error
bandwidth choices also need their labels retained. The inspected source
requires Stata version~16 semantics; this is a declared source requirement,
not a completed compatibility trial~\citep{exec13_ivrd_rdrobust}.

The \texttt{masspoints} options inspect repeated running-variable values.
The default adjustment enlarges preliminary bandwidths to include a
minimum number of distinct values, with ten as the documented default
and an explicit override. The source computes unique-score counts on
both sides for manual as well as automatically chosen bandwidths.
Its fuzzy branch rejects the fully degenerate constant treatment case
with no jump. That exact guard does not by itself establish coverage
when the first stage is merely small. Nor is robust bias correction
synonymous with weak-IV robustness. A matched trial must retain both
bandwidth uncertainty and weak local treatment uptake in its design
\citep{exec13_ivrd_rdrobust}.

The companion pinned \texttt{rddensity} help describes local-polynomial
density manipulation tests, bias-corrected and conventional statistics,
mass-point adjustment and complementary binomial testing. Its graphical
path requires \texttt{lpdensity}. Those are useful diagnostics, not a
certificate that applicants could not manipulate scores or that all
potential outcomes are continuous. A nondetectable density change can
coexist with another discontinuous policy; a density rejection can
reflect a measurement or design issue that needs investigation. The
proposal should retain test statistic, bandwidth, cutoff, support counts,
test family and qualification instead of converting a $p$-value into
a Boolean ``causal design verified'' field~\citep{exec13_ivrd_rddensity}.

Discrete scores show the remaining information problem. Suppose the
observed support is $\{-2,-1,1,2\}$, eligibility starts at zero, and
observed means are $80,90,130,140$. Let the treated response function
be $120+10s$. Two smooth untreated functions are
\[
 \mu_{0,\delta}(s)=100+10s+
             \frac{\delta}{2}(s+1)(s+2),\qquad \delta\in\{0,40\}.
\]
Both agree with the observed untreated values at $-2$ and $-1$.
Under sharp assignment the treated observations agree as well. Yet
the causal discontinuity at the unobserved cutoff is $20-\delta$,
which is $+20$ or $-20$. Repeating these same scores indefinitely
does not supply observations tending to zero. Smoothness without a
quantitative shape restriction does not choose between these laws.
A justified curvature bound, prespecified response model, local
randomization law or new score resolution is additional information.
Counting mass points is useful; it cannot create the missing support
\citep{exec13_ivrd_exact}.

The two upstream license files declare GPL version~3. The proposed
default is a reviewed adapter to a pinned distribution, with the
applicable license obligations assessed for the actual distribution
and use. Stata's host license is a separate issue. A port would need
to preserve the polynomial basis, kernel, bandwidth selection, bias
variance, cluster corrections, support accounting and missing-data
rules, with independent reference comparisons and an appropriate
license decision. Rewriting a few coefficient formulas would not
reproduce the package. Direct adoption inside a licensed Stata
workflow may be preferable for an existing user; a native R/Python
path may be preferable for another consumer, but cross-language
equivalence is an empirical obligation~\citep{exec13_ivrd_rdrobust,exec13_ivrd_rddensity}.

\subsection{The proposed product must compare complete decisions}

A useful IV/RD request carries the business decision, target population,
action and outcome units, decision and outcome clocks, cutoff rule,
running-variable precision, instrument origin, exclusion and compliance
assumptions, estimation sample and covariance design. The result retains
the target actually estimated, first-stage and support diagnostics,
every interval component, numerical search status, all failed requests
and the source/engine identities. A report must distinguish a
nonidentified effect, an unsupported operation, a scientifically
refused request and an unfinished numerical search.

The current local reference service can preserve generic supplied
moment and sensitivity sets, but it does not implement the acquired
commercial procedures or the pinned RD packages. The proposed
integration should begin with one licensed IV path and one pinned
RD adapter, each compared with a separate exact or small numerical
reference. The trial should include strong, weak and zero first stages;
valid and invalid exclusion; correlated and heteroskedastic observations;
coarse scores; manipulated cutoffs; alternative bandwidths; empty,
disconnected and whole-line sets; and grid exhaustion. It should
record all model fitting, conversion, failed-request and review costs.
These are proposed comparison cases, not completed vendor experiments.

The investment decision is whether those paths reduce errors and
repeated manual work enough to justify the integration. It is not
established by a longer list of estimator names. Stata's existing
weak-IV sets are a substantive baseline; SAS, RATS, EViews, SPSS
and JMP supply different supported entry points. The near-term
opportunity is to preserve their scientific meaning, expose their
limits, and connect a locally justified effect to a decision with
explicit costs. Measured savings and bank returns remain later
empirical questions~\citep{exec13_ivrd_survey,exec13_ivrd_exact}.
\label{exec13:commercial-iv-rd-end}
